% Standard preamble for literate programming documents
% Copy this file verbatim to doc/preamble.tex in new projects
%
% This preamble is shared across all literate programming projects.
% The main document (e.g., doc/main.tex) should % Standard preamble for literate programming documents
% Copy this file verbatim to doc/preamble.tex in new projects
%
% This preamble is shared across all literate programming projects.
% The main document (e.g., doc/main.tex) should % Standard preamble for literate programming documents
% Copy this file verbatim to doc/preamble.tex in new projects
%
% This preamble is shared across all literate programming projects.
% The main document (e.g., doc/main.tex) should % Standard preamble for literate programming documents
% Copy this file verbatim to doc/preamble.tex in new projects
%
% This preamble is shared across all literate programming projects.
% The main document (e.g., doc/main.tex) should \input{preamble} after
% \documentclass but before \begin{document}.

% Font setup.  Under pdfLaTeX these are REQUIRED: without [T1]{fontenc} a
% monospace [[code]] quote (or \texttt) cannot select a T1 typewriter shape and
% silently falls back to the proportional OT1 roman font -- i.e. monospace
% "disappears".  iftex guards them so this preamble stays safe when copied into
% a LuaLaTeX/XeLaTeX project (which already use Unicode fonts and need neither).
% Do NOT comment these out.  See latex-writing's references/unicode-and-fonts.md
% ("Symptom 2").
\usepackage{iftex}
\ifpdftex
  \usepackage[utf8]{inputenc}
  \usepackage[T1]{fontenc}
  \usepackage{lmodern}
\fi
\usepackage[swedish,british]{babel}
\usepackage{booktabs}

\usepackage{noweb}
\noweboptions{shift,breakcode,longxref,longchunks}

\usepackage[natbib,style=alphabetic,maxbibnames=99]{biblatex}
\addbibresource{bibliography.bib}

\usepackage[inline]{enumitem}
\setlist[enumerate]{label=(\arabic*)}

\usepackage[all]{foreign}
\renewcommand{\foreignfullfont}{}
\renewcommand{\foreignabbrfont}{}

%\usepackage{newclude}
\usepackage{import}

\usepackage[strict]{csquotes}
\usepackage[single]{acro}

\usepackage{subcaption}

\usepackage[noend]{algpseudocode}
\usepackage{xparse}

\let\email\texttt

%\usepackage[outputdir=ltxobj]{minted}
\usepackage{minted}
\setminted{autogobble}

\usepackage{pythontex}
\setpythontexoutputdir{.}
\setpythontexworkingdir{..}

\usepackage{fancyvrb}

\usepackage{amsmath}
\usepackage{amssymb}
\usepackage{mathtools}
\usepackage{amsthm}
\usepackage{thmtools}
\usepackage[unq]{unique}
\DeclareMathOperator{\powerset}{\mathcal{P}}

\usepackage[binary-units]{siunitx}

\usepackage{pgf}

\usepackage[colorlinks]{hyperref}
\usepackage[capitalize]{cleveref}
\crefalias{question}{item}
\crefname{question}{question}{questions}
\Crefname{question}{Question}{Questions}
\creflabelformat{question}{#2\textup{#1}#3}

% adapted from https://tex.stackexchange.com/a/30726/17418
\newcommand\chnote[1]{%
  \begingroup
  \renewcommand\thefootnote{}\footnote{#1}%
  \addtocounter{footnote}{-1}%
  \endgroup
}

\usepackage[marginparmargin=right]{didactic}

% Make noweb's hyperlinked identifiers safe in headings (chapter/section titles
% that contain a [[code]] quote).  With -index, such a quote weaves into a
% \nwlinkedidentq hyperlink; in a moving argument (the .toc, the running header,
% the PDF bookmark) hyperref then fails fatally (under newer TeX Live with
% errors like "Undefined control sequence \hyper@@link" / "\Hy@reserved@a", or a
% stack overflow from noweb's self-redefining \nwhyperreference).  Two fixes,
% both required, let the project keep [[...]] in headings instead of falling
% back to \texttt{...}:
%
% (1) Replace \nwhyperreference with noweb's robust hyperref form so it is
%     written literally into moving arguments instead of expanding (and
%     misfiring its terminating \global\let) there, yet still links in the body.
\DeclareRobustCommand\nwhyperreference[2]{\hyperlink{noweb.#1}{#2}}
%
% (2) Strip noweb's code-quote macros from PDF bookmark strings, so a [[code]]
%     quote in a heading yields clean bookmark text without hyperref warnings.
\pdfstringdefDisableCommands{%
  \def\nwlinkedidentq#1#2{#1}%
  \let\Tt\relax
  \let\nwendquote\relax
}

% Unicode character replacements for symbols not in Latin Modern fonts
% These appear in TUI code and documentation describing terminal output
\usepackage{newunicodechar}
\newunicodechar{✓}{\checkmark}
\newunicodechar{✗}{$\times$}
\newunicodechar{○}{\ensuremath{\circ}}
\newunicodechar{─}{\textemdash}
\newunicodechar{📝}{\texttt{[note]}}
\newunicodechar{🤖}{\texttt{[AI]}}
 after
% \documentclass but before \begin{document}.

% Font setup.  Under pdfLaTeX these are REQUIRED: without [T1]{fontenc} a
% monospace [[code]] quote (or \texttt) cannot select a T1 typewriter shape and
% silently falls back to the proportional OT1 roman font -- i.e. monospace
% "disappears".  iftex guards them so this preamble stays safe when copied into
% a LuaLaTeX/XeLaTeX project (which already use Unicode fonts and need neither).
% Do NOT comment these out.  See latex-writing's references/unicode-and-fonts.md
% ("Symptom 2").
\usepackage{iftex}
\ifpdftex
  \usepackage[utf8]{inputenc}
  \usepackage[T1]{fontenc}
  \usepackage{lmodern}
\fi
\usepackage[swedish,british]{babel}
\usepackage{booktabs}

\usepackage{noweb}
\noweboptions{shift,breakcode,longxref,longchunks}

\usepackage[natbib,style=alphabetic,maxbibnames=99]{biblatex}
\addbibresource{bibliography.bib}

\usepackage[inline]{enumitem}
\setlist[enumerate]{label=(\arabic*)}

\usepackage[all]{foreign}
\renewcommand{\foreignfullfont}{}
\renewcommand{\foreignabbrfont}{}

%\usepackage{newclude}
\usepackage{import}

\usepackage[strict]{csquotes}
\usepackage[single]{acro}

\usepackage{subcaption}

\usepackage[noend]{algpseudocode}
\usepackage{xparse}

\let\email\texttt

%\usepackage[outputdir=ltxobj]{minted}
\usepackage{minted}
\setminted{autogobble}

\usepackage{pythontex}
\setpythontexoutputdir{.}
\setpythontexworkingdir{..}

\usepackage{fancyvrb}

\usepackage{amsmath}
\usepackage{amssymb}
\usepackage{mathtools}
\usepackage{amsthm}
\usepackage{thmtools}
\usepackage[unq]{unique}
\DeclareMathOperator{\powerset}{\mathcal{P}}

\usepackage[binary-units]{siunitx}

\usepackage{pgf}

\usepackage[colorlinks]{hyperref}
\usepackage[capitalize]{cleveref}
\crefalias{question}{item}
\crefname{question}{question}{questions}
\Crefname{question}{Question}{Questions}
\creflabelformat{question}{#2\textup{#1}#3}

% adapted from https://tex.stackexchange.com/a/30726/17418
\newcommand\chnote[1]{%
  \begingroup
  \renewcommand\thefootnote{}\footnote{#1}%
  \addtocounter{footnote}{-1}%
  \endgroup
}

\usepackage[marginparmargin=right]{didactic}

% Make noweb's hyperlinked identifiers safe in headings (chapter/section titles
% that contain a [[code]] quote).  With -index, such a quote weaves into a
% \nwlinkedidentq hyperlink; in a moving argument (the .toc, the running header,
% the PDF bookmark) hyperref then fails fatally (under newer TeX Live with
% errors like "Undefined control sequence \hyper@@link" / "\Hy@reserved@a", or a
% stack overflow from noweb's self-redefining \nwhyperreference).  Two fixes,
% both required, let the project keep [[...]] in headings instead of falling
% back to \texttt{...}:
%
% (1) Replace \nwhyperreference with noweb's robust hyperref form so it is
%     written literally into moving arguments instead of expanding (and
%     misfiring its terminating \global\let) there, yet still links in the body.
\DeclareRobustCommand\nwhyperreference[2]{\hyperlink{noweb.#1}{#2}}
%
% (2) Strip noweb's code-quote macros from PDF bookmark strings, so a [[code]]
%     quote in a heading yields clean bookmark text without hyperref warnings.
\pdfstringdefDisableCommands{%
  \def\nwlinkedidentq#1#2{#1}%
  \let\Tt\relax
  \let\nwendquote\relax
}

% Unicode character replacements for symbols not in Latin Modern fonts
% These appear in TUI code and documentation describing terminal output
\usepackage{newunicodechar}
\newunicodechar{✓}{\checkmark}
\newunicodechar{✗}{$\times$}
\newunicodechar{○}{\ensuremath{\circ}}
\newunicodechar{─}{\textemdash}
\newunicodechar{📝}{\texttt{[note]}}
\newunicodechar{🤖}{\texttt{[AI]}}
 after
% \documentclass but before \begin{document}.

% Font setup.  Under pdfLaTeX these are REQUIRED: without [T1]{fontenc} a
% monospace [[code]] quote (or \texttt) cannot select a T1 typewriter shape and
% silently falls back to the proportional OT1 roman font -- i.e. monospace
% "disappears".  iftex guards them so this preamble stays safe when copied into
% a LuaLaTeX/XeLaTeX project (which already use Unicode fonts and need neither).
% Do NOT comment these out.  See latex-writing's references/unicode-and-fonts.md
% ("Symptom 2").
\usepackage{iftex}
\ifpdftex
  \usepackage[utf8]{inputenc}
  \usepackage[T1]{fontenc}
  \usepackage{lmodern}
\fi
\usepackage[swedish,british]{babel}
\usepackage{booktabs}

\usepackage{noweb}
\noweboptions{shift,breakcode,longxref,longchunks}

\usepackage[natbib,style=alphabetic,maxbibnames=99]{biblatex}
\addbibresource{bibliography.bib}

\usepackage[inline]{enumitem}
\setlist[enumerate]{label=(\arabic*)}

\usepackage[all]{foreign}
\renewcommand{\foreignfullfont}{}
\renewcommand{\foreignabbrfont}{}

%\usepackage{newclude}
\usepackage{import}

\usepackage[strict]{csquotes}
\usepackage[single]{acro}

\usepackage{subcaption}

\usepackage[noend]{algpseudocode}
\usepackage{xparse}

\let\email\texttt

%\usepackage[outputdir=ltxobj]{minted}
\usepackage{minted}
\setminted{autogobble}

\usepackage{pythontex}
\setpythontexoutputdir{.}
\setpythontexworkingdir{..}

\usepackage{fancyvrb}

\usepackage{amsmath}
\usepackage{amssymb}
\usepackage{mathtools}
\usepackage{amsthm}
\usepackage{thmtools}
\usepackage[unq]{unique}
\DeclareMathOperator{\powerset}{\mathcal{P}}

\usepackage[binary-units]{siunitx}

\usepackage{pgf}

\usepackage[colorlinks]{hyperref}
\usepackage[capitalize]{cleveref}
\crefalias{question}{item}
\crefname{question}{question}{questions}
\Crefname{question}{Question}{Questions}
\creflabelformat{question}{#2\textup{#1}#3}

% adapted from https://tex.stackexchange.com/a/30726/17418
\newcommand\chnote[1]{%
  \begingroup
  \renewcommand\thefootnote{}\footnote{#1}%
  \addtocounter{footnote}{-1}%
  \endgroup
}

\usepackage[marginparmargin=right]{didactic}

% Make noweb's hyperlinked identifiers safe in headings (chapter/section titles
% that contain a [[code]] quote).  With -index, such a quote weaves into a
% \nwlinkedidentq hyperlink; in a moving argument (the .toc, the running header,
% the PDF bookmark) hyperref then fails fatally (under newer TeX Live with
% errors like "Undefined control sequence \hyper@@link" / "\Hy@reserved@a", or a
% stack overflow from noweb's self-redefining \nwhyperreference).  Two fixes,
% both required, let the project keep [[...]] in headings instead of falling
% back to \texttt{...}:
%
% (1) Replace \nwhyperreference with noweb's robust hyperref form so it is
%     written literally into moving arguments instead of expanding (and
%     misfiring its terminating \global\let) there, yet still links in the body.
\DeclareRobustCommand\nwhyperreference[2]{\hyperlink{noweb.#1}{#2}}
%
% (2) Strip noweb's code-quote macros from PDF bookmark strings, so a [[code]]
%     quote in a heading yields clean bookmark text without hyperref warnings.
\pdfstringdefDisableCommands{%
  \def\nwlinkedidentq#1#2{#1}%
  \let\Tt\relax
  \let\nwendquote\relax
}

% Unicode character replacements for symbols not in Latin Modern fonts
% These appear in TUI code and documentation describing terminal output
\usepackage{newunicodechar}
\newunicodechar{✓}{\checkmark}
\newunicodechar{✗}{$\times$}
\newunicodechar{○}{\ensuremath{\circ}}
\newunicodechar{─}{\textemdash}
\newunicodechar{📝}{\texttt{[note]}}
\newunicodechar{🤖}{\texttt{[AI]}}
 after
% \documentclass but before \begin{document}.

% Font setup.  Under pdfLaTeX these are REQUIRED: without [T1]{fontenc} a
% monospace [[code]] quote (or \texttt) cannot select a T1 typewriter shape and
% silently falls back to the proportional OT1 roman font -- i.e. monospace
% "disappears".  iftex guards them so this preamble stays safe when copied into
% a LuaLaTeX/XeLaTeX project (which already use Unicode fonts and need neither).
% Do NOT comment these out.  See latex-writing's references/unicode-and-fonts.md
% ("Symptom 2").
\usepackage{iftex}
\ifpdftex
  \usepackage[utf8]{inputenc}
  \usepackage[T1]{fontenc}
  \usepackage{lmodern}
\fi
\usepackage[swedish,british]{babel}
\usepackage{booktabs}

\usepackage{noweb}
\noweboptions{shift,breakcode,longxref,longchunks}

\usepackage[natbib,style=alphabetic,maxbibnames=99]{biblatex}
\addbibresource{bibliography.bib}

\usepackage[inline]{enumitem}
\setlist[enumerate]{label=(\arabic*)}

\usepackage[all]{foreign}
\renewcommand{\foreignfullfont}{}
\renewcommand{\foreignabbrfont}{}

%\usepackage{newclude}
\usepackage{import}

\usepackage[strict]{csquotes}
\usepackage[single]{acro}

\usepackage{subcaption}

\usepackage[noend]{algpseudocode}
\usepackage{xparse}

\let\email\texttt

%\usepackage[outputdir=ltxobj]{minted}
\usepackage{minted}
\setminted{autogobble}

\usepackage{pythontex}
\setpythontexoutputdir{.}
\setpythontexworkingdir{..}

\usepackage{fancyvrb}

\usepackage{amsmath}
\usepackage{amssymb}
\usepackage{mathtools}
\usepackage{amsthm}
\usepackage{thmtools}
\usepackage[unq]{unique}
\DeclareMathOperator{\powerset}{\mathcal{P}}

\usepackage[binary-units]{siunitx}

\usepackage{pgf}

\usepackage[colorlinks]{hyperref}
\usepackage[capitalize]{cleveref}
\crefalias{question}{item}
\crefname{question}{question}{questions}
\Crefname{question}{Question}{Questions}
\creflabelformat{question}{#2\textup{#1}#3}

% adapted from https://tex.stackexchange.com/a/30726/17418
\newcommand\chnote[1]{%
  \begingroup
  \renewcommand\thefootnote{}\footnote{#1}%
  \addtocounter{footnote}{-1}%
  \endgroup
}

\usepackage[marginparmargin=right]{didactic}

% Make noweb's hyperlinked identifiers safe in headings (chapter/section titles
% that contain a [[code]] quote).  With -index, such a quote weaves into a
% \nwlinkedidentq hyperlink; in a moving argument (the .toc, the running header,
% the PDF bookmark) hyperref then fails fatally (under newer TeX Live with
% errors like "Undefined control sequence \hyper@@link" / "\Hy@reserved@a", or a
% stack overflow from noweb's self-redefining \nwhyperreference).  Two fixes,
% both required, let the project keep [[...]] in headings instead of falling
% back to \texttt{...}:
%
% (1) Replace \nwhyperreference with noweb's robust hyperref form so it is
%     written literally into moving arguments instead of expanding (and
%     misfiring its terminating \global\let) there, yet still links in the body.
\DeclareRobustCommand\nwhyperreference[2]{\hyperlink{noweb.#1}{#2}}
%
% (2) Strip noweb's code-quote macros from PDF bookmark strings, so a [[code]]
%     quote in a heading yields clean bookmark text without hyperref warnings.
\pdfstringdefDisableCommands{%
  \def\nwlinkedidentq#1#2{#1}%
  \let\Tt\relax
  \let\nwendquote\relax
}

% Unicode character replacements for symbols not in Latin Modern fonts
% These appear in TUI code and documentation describing terminal output
\usepackage{newunicodechar}
\newunicodechar{✓}{\checkmark}
\newunicodechar{✗}{$\times$}
\newunicodechar{○}{\ensuremath{\circ}}
\newunicodechar{─}{\textemdash}
\newunicodechar{📝}{\texttt{[note]}}
\newunicodechar{🤖}{\texttt{[AI]}}
